\chapter{2015年四川卷（理）}

\section{单选题}

\begin{problem}
设集合 \(A = \{x \mid (x + 1)(x - 2) < 0\}\)，集合 \(B = \{x \mid 1 < x < 3\}\)，则 \(A \cup B =\)（\quad）

\choices
    {\(\{x\mid -1 < x < 3\}\)}
    {\(\{x\mid -1 < x < 1\}\)}
    {\(\{x\mid 1 < x < 2\}\)}
    {\(\{x\mid 2 < x < 3\}\)}
\end{problem}

\begin{problem}
设 \(i\) 是虚数单位，则复数 \(i^{3}-\frac{2}{i}= \)（\quad）

\choices
    {-i}
    {-3i}
    {i}
    {3i}
\end{problem}

\begin{problem}
执行如图所示的程序框图，输出 \(S\) 的值为（\quad）

\texfigure[max width=6.0cm]{img_repaint/2015/sichuan_science/q3_fig1.tex}

\choices
    {\(-\frac{\sqrt{3}}{2}\)}
    {\(\frac{\sqrt{3}}{2}\)}
    {\(-\frac{1}{2}\)}
    {\(\frac{1}{2}\)}
\end{problem}

\begin{problem}
下列函数中，最小正周期为 \(\pi\) 且图象关于原点对称的函数是（\quad）

\choices
    {\( y = \cos \left(2x + \frac{\pi}{2}\right) \)}
    {\( y = \sin \left(2x + \frac{\pi}{3}\right) \)}
    {\( y = \sin 2x + \cos 2x \)}
    {\( y = \sin x + \cos x \)}
\end{problem}

\begin{problem}
过双曲线 \(x^{2} - \frac{y^{2}}{3} = 1\) 的右焦点且与 \(x\) 轴垂直的直线，交该双曲线的两条渐近线于 \(A\)，\(B\) 两点，\(|AB| =\)（\quad）

\choices
    {\(\frac{4\sqrt{3}}{3}\)}
    {\(2 \sqrt{3}\)}
    {6}
    {\(4 \sqrt{3}\)}
\end{problem}

\begin{problem}
用数字 0， 1， 2， 3， 4， 5 组成没有重复数字的五位数，其中比 40000 大的偶数共有（\quad）

\choices
    {144 个}
    {120 个}
    {96 个}
    {72 个}
\end{problem}

\begin{problem}
设四边形 \(ABCD\) 为平行四边形，\(\left|\overrightarrow{AB}\right| = 6\)，\(\left|\overrightarrow{AD}\right| = 4\). 若点 \(M\)，\(N\) 满足 \(\overrightarrow{BM} = 3\overrightarrow{MC}\)，\(\overrightarrow{DN} = 2\overrightarrow{NC}\)，则 \(\overrightarrow{AM} \cdot \overrightarrow{NM} =\)（\quad）

\choices
    {20}
    {15}
    {9}
    {6}
\end{problem}

\begin{problem}
设 \(a\)，\(b\) 都是不等于 1 的正数，则 \(3^{a} > 3^{b} > 3\) 是“\(\log_{a} 3 < \log_{b} 3\)”的（\quad）

\choices
    {充要条件}
    {充分不必要条件}
    {必要不充分条件}
    {既不充分也不必要条件}
\end{problem}

\begin{problem}
如果函数 \( f(x) = \frac{1}{2} (m - 2)x^2 + (n - 8)x + 1  \) （\(m \geqslant 0\)，\(n \geqslant 0\)）在区间 \(\left[\frac{1}{2}, 2\right]\) 上单调递减，那么 \( mn \) 的最大值为（\quad）

\choices
    {16}
    {18}
    {25}
    {\(\frac{81}{2}\)}
\end{problem}

\begin{problem}
设直线 \(l\) 与抛物线 \(y^{2} = 4x\) 相交于 \(A\)，\(B\) 两点，与圆 \((x - 5)^{2} + y^{2} = r^{2}\)（\(r > 0\)）相切于点 \(M\)，且 \(M\) 为线段 \(AB\) 的中点. 若这样的直线 \(l\) 恰有 4 条，则 \(r\) 的取值范围是（\quad）

\choices
    {(1,3)}
    {(1,4)}
    {(2,3)}
    {(2,4)}
\end{problem}

\section{填空题}

\begin{problem}
在 \((2x - 1)^{5}\) 的展开式中，含 \(x^{2}\) 项的系数是\fillinblank{}. （用数字填写答案）
\end{problem}

\begin{problem}
\(\sin 15^{\circ} + \sin 75^{\circ}\) 的值是\fillinblank{}.
\end{problem}

\begin{problem}
某食品的保鲜时间 \(y\)（单位：小时）与储藏温度 \(x\)（单位：\(^{\circ}\mathrm{C}\)）满足函数关系 \(y = \e^{kx + b}\)（\(\e = 2.718 \cdots\) 为自然对数的底数，\(k\)，\(b\) 为常数）. 若该食品在 \(0^{\circ}\mathrm{C}\) 的保鲜时间是 192 小时，在 \(22^{\circ}\mathrm{C}\) 的保鲜时间是 48 小时，则该食品在 \(33^{\circ}\mathrm{C}\) 的保鲜时间是\fillinblank{} 小时.
\end{problem}

\begin{problem}
如图，四边形 \(ABCD\) 和 \(ADPQ\) 均为正方形，它们所在的平面互相垂直，动点 \(M\) 在线段 \(PQ\) 上，\(E\)，\(F\) 分别为 \(AB\)，\(BC\) 的中点. 设异面直线 \(EM\) 与 \(AF\) 所成的角为 \(\theta\)，则 \(\cos \theta\) 的最大值为\fillinblank{}.

\bitmapfigure[max width=0.36\linewidth]{img/2015/sichuan_science/q14_fig1.png}
\end{problem}

\begin{problem}
已知函数 \(f(x) = 2^x\)，\(g(x) = x^2 + ax\)（其中 \(a \in \R\)）.对于不相等的实数 \(x_1\)，\(x_2\)，设 \(m = \frac{f(x_1) - f(x_2)}{x_1 - x_2}\)，\(n = \frac{g(x_1) - g(x_2)}{x_1 - x_2}\)，现有如下命题：
\begin{circlelist}
\item 对于任意不相等的实数 \(x_{1}\)，\(x_{2}\)，都有 \(m > 0\)；
\item 对于任意的 \(a\) 及任意不相等的实数 \(x_{1}\)，\(x_{2}\)，都有 \(n > 0\)；
\item 对于任意的 \(a\)，存在不相等的实数 \(x_{1}\)，\(x_{2}\)，使得 \(m = n\)；
\item 对于任意的 \(a\)，存在不相等的实数 \(x_{1}\)，\(x_{2}\)，使得 \(m = -n\).
\end{circlelist}
其中的真命题有\fillinblank{}. （写出所有真命题的序号）
\end{problem}

\section{解答题}

\begin{problem}
设数列 \(\{a_{n}\} (n = 1,2,3,\dots)\) 的前 \(n\) 项和 \(S_{n}\) 满足 \(S_{n} = 2a_{n} - a_{1}\) ，且 \(a_1\)，\(a_2 + 1\)，\(a_3\) 成等差数列.
\begin{enumerate}
\item 求数列 \( \{a_{n}\} \) 的通项公式；
\item 记数列 \( \left\{\frac{1}{a_{n}}\right\} \) 的前 \(n\) 项和为 \( T_{n} \)，求使得 \( \left|T_{n} - 1\right| < \frac{1}{1000} \) 成立的 \(n\) 的最小值.
\end{enumerate}
\end{problem}

\begin{problem}
某市 A, B 两所中学的学生组队参加辩论赛，A 中学推荐了 3 名男生，2 名女生，B 中学推荐了 3 名男生，4 名女生. 两校所推荐的学生一起参加集训. 由于集训后队员水平相当，从参加集训的男生中随机抽取 3 人，女生中随机抽取 3 人组成代表队.
\begin{enumerate}
\item 求 A 中学至少有 1 名学生入选代表队的概率；
\item 某场比赛前，从代表队的 6 名队员中随机抽取 4 人参赛，设 \(X\) 表示参赛的男生人数，求 \(X\) 的分布列和数学期望.
\end{enumerate}
\end{problem}

\begin{problem}
一个正方体的平面展开图及该正方体的直观图的示意图如图所示. 在正方体中，设 \(BC\) 的中点为 \(M\)，\(GH\) 的中点为 \(N\).
\begin{enumerate}
\item 请将字母 \(F\)，\(G\)，\(H\) 标记在正方体相应的顶点处（不需说明理由）；
\item 证明：直线 \(MN\parallel\) 平面 \(BDH\)；
\item 求二面角 \(A - EG - M\) 的余弦值.
\end{enumerate}
\end{problem}

\begin{problem}
如图，\(A\)，\(B\)，\(C\)，\(D\) 为平面四边形 \(ABCD\) 的四个内角.
\begin{enumerate}
\item 证明：\(\tan \frac{A}{2} = \frac{1 - \cos A}{\sin A}\)；
\item 若 \(A + C = 180^{\circ}\)，\(AB = 6\)，\(BC = 3\)，\(CD = 4\)，\(AD = 5\)， 求 \(\tan \frac{A}{2} + \tan \frac{B}{2} + \tan \frac{C}{2} + \tan \frac{D}{2}\) 的值.
\end{enumerate}

\bitmapfigure[max width=0.38\linewidth]{img/2015/sichuan_science/q19_fig1.png}
\end{problem}

\begin{problem}
如图，椭圆 \(E: \frac{x^2}{a^2} + \frac{y^2}{b^2} = 1 \) （\(a > b > 0\)）的离心率是 \(\frac{\sqrt{2}}{2}\)，过点 \(P(0,1)\) 的动直线 \(l\) 与椭圆相交于 \(A\)，\(B\) 两点. 当直线 \(l\) 平行于 \(x\) 轴时，直线 \(l\) 被椭圆 \(E\) 截得的线段长为 \(2\sqrt{2}\).

\begin{enumerate}
\item 求椭圆 \(E\) 的方程；
\item 在平面直角坐标系 \(xOy\) 中，是否存在与点 \(P\) 不同的定点 \(Q\)，使得 \(\frac{|QA|}{|QB|} = \frac{|PA|}{|PB|}\) 恒成立? 若存在，求出点 \(Q\) 的坐标；若不存在，请说明理由.
\end{enumerate}
\bitmapfigure[max width=0.42\linewidth]{img/2015/sichuan_science/q20_fig1.png}
\end{problem}

\begin{problem}
已知函数 \( f(x) = -2(x + a)\ln x + x^2 - 2ax - 2a^2 + a \)，其中 \( a > 0 \).
\begin{enumerate}
\item 设 \(g(x)\) 是 \(f(x)\) 的导函数，讨论 \(g(x)\) 的单调性；
\item 证明：存在 \(a \in (0,1)\)，使得 \(f(x) \geqslant 0\) 在区间 \((1, +\infty)\) 内恒成立，且 \(f(x) = 0\) 在区间 \((1, +\infty)\) 内有唯一解.
\end{enumerate}
\end{problem}
